\documentclass[11pt]{article}
\usepackage[margin=1in]{geometry}
\usepackage{amsmath,amssymb,amsthm}
\newtheorem{proposition}{Proposition}
\newcommand{\Id}{\mathbb I}
\begin{document}
\begin{center}\large A quantum-valid extreme curve with two deterministic zeros\end{center}

Let $A_0,A_1,B_0,B_1$ be binary observables, with the two parties
commuting, and put
\[
 S=A_0B_0-A_0B_1+A_1B_0+A_1B_1,\qquad H=\Id-S/2.
\]
Write $p_{ab|xy}$ for the joint event probability, with $a,b\in\{+,-\}$.
Our canonical slack is
\begin{align*}
 f={}&dH+4\lambda p_{++|01}+4n p_{+-|00}+4m p_{+-|11}
       +4C p_{-+|10}\\
    &+4A p_{-+|00}+4D p_{-+|11}
       +4\alpha p_{--|01}+4\beta p_{+-|10}.
\end{align*}
Define the canonical valid cone as the set of its nonnegative coefficient
vectors for which $f\geq0$ on every quantum behaviour.

\begin{proposition}
For $0<z\leq2/5$, the source $d=1$, $\lambda=\beta=0$,
\begin{align*}
 n&=\frac{2-5z}{8(1+2z)},&
 m&=\frac{2+z}{8(1+2z)},&
 C&=\frac{2+z}{2(1+2z)},\\
 A&=\frac{3(2+z)}{8(1+2z)},&
 D&=\frac{3(2-z)}{8(1+2z)},&
 \alpha&=\frac{3z^2}{2(2-z)(1+2z)}
\end{align*}
is quantum-valid and generates an extreme ray of the canonical valid cone.
The assertion includes $z=2/5$, where $n=0$.
\end{proposition}

\begin{proof}
We prove validity by an explicit determinant decomposition.
Multiplication by $K=8(2-z)(1+2z)>0$ gives the polynomial source
\[
\begin{split}
 (d,\lambda,n,m,C,A,D,\alpha,\beta)
 =\bigl(&8(2-z)(1+2z),0,(2-z)(2-5z),\,4-z^2,\\
       &4(4-z^2),\,3(4-z^2),\,3(2-z)^2,\,12z^2,0\bigr).
\end{split}                                                    \tag{1}
\]
All entries are nonnegative, and $d>0$. Use real two-qubit coordinates
\[
 A_0=B_0=Z,\quad
 A_1=\frac{(1-u^2)Z+2uX}{1+u^2},\quad
 B_1=\frac{(1-v^2)Z+2vX}{1+v^2}.
\]
For a real state $\psi=(W,X_s,Y,Z_s)$, put $D_u=1+u^2$, $D_v=1+v^2$.
The entire cleared slack for the source (1) is the quadratic form
\begin{align*}
 F={}&4nD_uD_vX_s^2+4m(vW-X_s+uvY-uZ_s)^2
       +4CD_v(uW-Y)^2+4AD_uD_vY^2\\
 &+(4D+2d)(uW+uvX_s-Y-vZ_s)^2
 +(4\alpha+2d)D_u(vY-Z_s)^2\\
 &+2dD_v(X_s+uZ_s)^2-2dD_uD_vZ_s^2.
                                                               \tag{2}
\end{align*}
Thus $F=D_uD_v\|\psi\|^2 f(q)$, and we write $P_z(u,v)$ for its
symmetric state matrix. Its $Z_s=0$ principal block is positive definite
when $uv\ne0$: the terms proportional to $d$ alone give
\[
 2d\bigl[(uW+uvX_s-Y)^2+D_uv^2Y^2+D_vX_s^2\bigr].       \tag{3}
\]
Consequently $P_z$ has three positive eigenvalues; the sign of its fourth
eigenvalue is the sign of its determinant.

Here the determinant has a particularly short exact form. Set
$U=u^2$, $V=v^2$, $G_z=432(2-z)^3(z+2)^3$. Direct expansion of (2) gives
\[
 \det P_z=16G_zD_u^2D_v^2u^2v^2 R_z(U,V),                \tag{4}
\]
where
\begin{align*}
R_z(U,V)={}&14z^2-96z-24+\frac{16z^2}{U}
 +(6-z)(2-z)\frac VU+\frac{48z^2}{V}\\
 &+16(2-z)V+3(z+2)(2-5z)\frac UV+16(z+2)U.
                                                               \tag{5}
\end{align*}
For completeness, (4) can be checked without a four-dimensional
determinant expansion by comparing the six outer coefficients in (5)
and the constant coefficient after setting $u=v=1$. The accompanying
exact checker also expands the full polynomial determinant in $(u,v,z)$.

Put $x=U/z$, $y=V/z$. The whole Laurent polynomial is
\begin{align*}
R_z(zx,zy)={}&16z(x+x^{-1}-2)
 +16z(2-z)(y+y^{-1}-2)\\
 &+3(z+2)(2-5z)(x/y+y/x-2)\\
 &+16z(z+1)(y/x+x+y^{-1}-3).                              \tag{6}
\end{align*}
Every coefficient is nonnegative on the stated interval. The first three
brackets are nonnegative reciprocal-pair expressions. The last is
nonnegative by three-term AM--GM, because $(y/x)x y^{-1}=1$.
This also exhibits balance and a zero at $U=V=z$ directly.
Equations (3)--(6) give $P_z(u,v)\succeq0$ for $uv\ne0$.
Continuity of $P_z/(D_uD_v)$ gives positivity on the axes and at projective
half-angle endpoints. Any two binary observables decompose into blocks of dimension
at most two; in each block a local basis gives the displayed real
half-angle form. Positivity of its real state matrix also gives positivity
on complex state vectors. Hence this proves quantum validity for all
states and binary measurements.

It remains to establish extremality. Return to the unscaled source
and let $c=\sqrt z$. At $u=v=c$ its Bell operator has the kernel vector
\[
 \psi_c=(1,-c/2,c/2,1)^T.                                 \tag{7}
\]
Substitution into (2), or into the Bell operator before clearing, verifies
the four state-kernel equations and the two angular first-derivative
equations. To check their independence, regard them as linear equations
in the seven active coordinates $(d,n,m,C,A,D,\alpha)$, holding (7)
fixed while differentiating the measurement angles. More precisely, let
$\mathcal B(u,v)$ be the uncleared Bell operator $f$, and define $M(c)$
as the coefficient matrix, in that ordered source basis, of
\[
 \left.\begin{pmatrix}
 \mathcal B\psi_c\\
 \psi_c^T(\partial_u\mathcal B)\psi_c\\
 \psi_c^T(\partial_v\mathcal B)\psi_c
 \end{pmatrix}\right|_{u=v=c}.
\]
Here $\psi_c$ is unnormalised, the first four rows are state-kernel
entries rather than doubled quadratic gradients, and the final two rows
are scalar angular derivatives. Then
\[
 \det M(c)_{[:,\,(d,n,m,C,A,D)]}
       =-\frac{768c^{10}}{(1+c^2)^5}\ne0.                 \tag{8}
\]
Thus their solution space is exactly the one-dimensional span of our
source. Now split it as a sum of two canonical valid sources.
Nonnegativity preserves $\lambda=\beta=0$ in both summands. Each summand
has zero value at (7), since it is valid and their sum vanishes there.
Its positive semidefinite Bell operator therefore annihilates $\psi_c$;
its two angular derivatives vanish at this interior minimum. By (8),
each summand is proportional to the original source. The proportion is
nonnegative because its $d$ coordinate is nonnegative. This proves that
the ray is extreme, including the additional face $n=0$.
\end{proof}
\end{document}
